\documentclass{scrreprt}
\usepackage{listings}
\usepackage{underscore}
\usepackage{graphicx}
\usepackage[bookmarks=true]{hyperref}
\usepackage[utf8]{inputenc}
\usepackage[english]{babel}
\hypersetup{
    bookmarks=false,    % show bookmarks bar?
    pdftitle={Software Requirement Specification},    % title
    pdfauthor={Jean-Philippe Eisenbarth},                     % author
    pdfsubject={TeX and LaTeX},                        % subject of the document
    pdfkeywords={TeX, LaTeX, graphics, images}, % list of keywords
    colorlinks=true,       % false: boxed links; true: colored links
    linkcolor=blue,       % color of internal links
    citecolor=black,       % color of links to bibliography
    filecolor=black,        % color of file links
    urlcolor=purple,        % color of external links
    linktoc=page            % only page is linked
}%
\def\myversion{1.0 }
\date{}
%\title{%

%}
\usepackage{hyperref}
\begin{document}

\begin{flushright}
    \rule{16cm}{5pt}\vskip1cm
    \begin{bfseries}
        \Huge{SOFTWARE REQUIREMENTS\\ SPECIFICATION}\\
        \vspace{1.5cm}
        for\\
        \vspace{1.5cm}
        Vending Machine Restocking Assistant\\
        \vspace{1.5cm}
        \LARGE{Version \myversion}\\
        \vspace{1.5cm}
        Prepared by:\\
        Joel Sarmiento\\
        Devin Schupbach\\
        Caden Sprague\\
        \vspace{1.5cm}
        July 16, 2026\\
    \end{bfseries}
\end{flushright}

\tableofcontents

\chapter{Introduction}

\section{Purpose -- JOEL}
The purpose of this Software Requirements Specification is to define the functional and non-functional requirements for the Vending Machine Restocking Assistant, a mobile-first application and supporting backend service designed to streamline the Micromart restocking process. This document describes what the system shall do, the constraints under which it must operate, and the quality attributes it must satisfy.

\section{Intended Audience -- JOEL}
This document is intended for the development team, who will use it as the primary reference for design, implementation, and testing; representatives of the vending machine stocking company, who may use it to verify that the proposed system aligns with their operational needs; and testers, who will use the functional and non-functional requirements to verify system behavior.

\section{Product Scope -- JOEL}
The Vending Machine Restocking Assistant is a mobile application and backend service that allows vending machine stockers to update Micromart inventory using simple spoken or typed commands. The system has the following functions:
\begin{itemize}
    \item Load the active site and restocking session information at the start of a session
    \item Fetch and index the machine's planogram from Micromart
    \item Accept spoken or typed restock commands from the stocker
    \item Match product names in commands to Micromart inventory records
    \item Prompt the stocker to clarify when a command matches multiple products
    \item Submit restock correction events to the Micromart system
    \item Display success or failure feedback to the stocker
    \item Log all attempted updates for debugging purposes
\end{itemize}

\section{Definitions, Acronyms, and Abbreviations -- JOEL}
\begin{tabular}{|p{4cm}|p{10cm}|}
\hline
\textbf{Term} & \textbf{Definition} \\
\hline
Stocker & A field employee responsible for physically restocking a vending machine \\
\hline
Planogram & A schematic provided by Micromart that defines the products and slots in a vending machine \\
\hline
Restock event & An API action submitted to Micromart that updates the recorded inventory quantity for a product \\
\hline
Session cookie & An authentication token used to authorize requests to the Micromart system \\
\hline
Micromart & The third-party vending machine management platform used by the stocking company \\
\hline
API & Application Programming Interface \\
\hline
Backend & The server-side service responsible for command parsing, product matching, and Micromart communication \\
\hline
\end{tabular}

\section{References -- JOEL}
\begin{tabular}{|p{4cm}|p{10cm}|}
\hline
\textbf{Reference} & \textbf{Source} \\
\hline
IEEE Std 830-1998 & IEEE Recommended Practice for Software Requirements Specifications \\
\hline
Micromart API & Internal Micromart web platform used for inventory management \\
\hline
\end{tabular}
\section{Overview -- JOEL}
This document is organized using the outline formate laid out by IEEE. It contains overall description of the system, functional requirements, non-functional requirements, and design constraints.

\chapter{Overall Description}

\section{Product Perspective -- CADEN}
The Vending Machine Restocking Assistant is a new system consisting of a mobile-first app and a backend service. It is not a standalone product but an intermediary that interacts with the existing, third-party Micromart platform: the backend communicates with Micromart to fetch planogram data and submit restock correction events on the stocker's behalf. The project team is not responsible for building or modifying the Micromart platform itself, the vending machine hardware, or the company's existing website. The stocking company currently stocks more than one type of vending machine system, including 365 and Micromart machines, but this system will focus first on Micromart restocking.

\section{User Classes and Characteristics -- CADEN}
The primary user class is the stocker, a field employee who physically restocks vending machines and needs a faster way to update inventory than interacting with the existing Micromart website. For the initial version, the system will support one user and one active restocking session at a time.

\section{Product Functions -- CADEN}
The system will support the following high-level functions:
\begin{itemize}
    \item Load the active site and restock session information at the beginning of a restocking session
    \item Fetch the machine's planogram from Micromart and build a searchable inventory index from it
    \item Receive spoken or typed restock commands from the stocker, such as ``set Fairlife to 5''
    \item Parse restock commands and match product names to Micromart inventory records
    \item Prompt the stocker to choose the correct item when a command matches more than one product
    \item Handle unclear or invalid commands
    \item Submit restock correction events to the Micromart system
    \item Show the stocker whether an update succeeded or failed
    \item Log attempted updates for debugging and safety
\end{itemize}

\section{Operating Environment -- CADEN}
The system consists of a mobile app used by the stocker in the field and a backend service that communicates with the Micromart system. The mobile app is responsible for capturing user input, displaying results, and handling product clarification when needed. The backend is responsible for Micromart communication, product matching logic, command parsing, API error handling, and restock event submission.

\section{Design and Implementation Constraints -- CADEN}
For the initial version, the system will support only one user and one active restocking session at a time. Authentication with Micromart will be simplified for the project MVP by using a manually supplied session cookie, but the backend will be designed so that a real login and two-factor authentication flow could be added later. The project team is not responsible for building or modifying the Micromart platform itself, the vending machine hardware, or the company's existing website.

\section{Assumptions and Dependencies -- CADEN}
The system assumes a valid, manually supplied Micromart session cookie is available for authentication in the MVP, rather than a full login and two-factor authentication flow. The system depends on the Micromart platform for planogram data and for accepting submitted restock correction events, and the backend's ability to match products and submit updates depends on the accuracy of the planogram fetched at the start of each session.

\chapter{System Features}

\section{Restocking Session Initialization -- CADEN}
At the beginning of a restocking session, the system loads the active site and restock session information. The backend then fetches the active machine's planogram from Micromart and builds a searchable inventory index from it, so that product names given in later commands can be matched against the actual products in the machine.

\textbf{Priority:} Essential

\textbf{Functional Requirements:}
\begin{itemize}
    \item R1.1 The system shall load the active site and restock session information at the start of a restocking session.
    \item R1.2 The system shall fetch the active machine's planogram from Micromart at the start of a restocking session.
    \item R1.3 The system shall build a searchable inventory index from the fetched planogram before accepting restock commands.
\end{itemize}

\section{Voice and Text Command Input -- JOEL}
The system accepts restock commands from the stocker as either spoken 
voice input or typed text. Voice input is converted to text by the 
mobile device before being sent to the backend. A valid command 
consists of a product name and a target quantity, for example 
``set Fairlife to 5.'' The backend parses the command string and 
extracts these two values. If the command cannot be parsed into a 
recognizable product name and quantity, the system returns an error 
message to the app. The system shall return a result within 3 seconds 
of receiving a command.

\textbf{Priority:} Essential

\textbf{Functional Requirements:}
\begin{itemize}
    \item R2.1 The system shall accept restock commands as typed text input.
    \item R2.2 The system shall accept restock commands as spoken voice 
    input converted to text by the mobile device.
    \item R2.3 The system shall parse a valid command into a product name 
    and a non-negative integer quantity.
    \item R2.4 The system shall return an invalid command error to the app 
    if the input cannot be parsed.
    \item R2.5 The system shall return a result to the app within 3 seconds 
    of receiving a command.
\end{itemize}

\section{Product Name Matching -- JOEL}
At the start of each restocking session the backend builds a searchable 
product index from the machine's planogram. When a command is received, 
the system searches the index for the closest matching product using 
fuzzy name matching to account for spelling variations and abbreviations. 
The system returns one of three results: a single matched product with 
its Micromart inventory identifier, a list of candidates if multiple 
products match, or a not found result if no match exists.

\textbf{Priority:} Essential

\textbf{Functional Requirements:}
\begin{itemize}
    \item R3.1 The system shall build a searchable product index from the 
    planogram at the start of each restocking session.
    \item R3.2 The system shall use fuzzy string matching to compare parsed 
    product names against the index.
    \item R3.3 The system shall return a matched product and its Micromart 
    inventory identifier when exactly one match is found.
    \item R3.4 The system shall return a list of candidates when a command 
    matches more than one product.
    \item R3.5 The system shall return a not found result when no product 
    in the index matches the parsed name.
\end{itemize}

\section{Product Disambiguation -- JOEL}
When a restock command matches more than one product in the planogram 
index, the system sends the list of candidates to the mobile app and 
waits for the stocker to select the correct item before submitting any 
update. The system shall not submit a restock correction event until a 
single product has been confirmed. The stocker may also cancel the 
command entirely, in which case no update is submitted.

\textbf{Priority:} Essential

\textbf{Functional Requirements:}
\begin{itemize}
    \item R4.1 The system shall send all matching candidate products to 
    the app when a command is ambiguous.
    \item R4.2 The system shall not submit a restock correction event until 
    the stocker selects a single product from the candidates.
    \item R4.3 The system shall allow the stocker to cancel a disambiguation 
    prompt without submitting any update.
    \item R4.4 The system shall use the originally parsed quantity when 
    submitting after a product has been selected.
\end{itemize}

\section{Restock Correction Submission -- CADEN}
Once a restock command has resolved to a single matched product, whether directly or after the stocker resolves a disambiguation prompt, the backend submits the corresponding restock correction event to the Micromart system. The app then shows the stocker whether the update succeeded or failed.

\textbf{Priority:} Essential

\textbf{Functional Requirements:}
\begin{itemize}
    \item R5.1 The system shall submit a restock correction event to Micromart once a command has resolved to a single matched product and quantity.
    \item R5.2 The system shall use the Micromart inventory identifier of the matched product when submitting the restock correction event.
    \item R5.3 The system shall report to the app whether the submitted restock correction event succeeded or failed.
\end{itemize}

\section{Update Result Display -- DEVIN}
Once a restock correction event has been submitted to Micromart, or once the system has determined that a command could not be completed, the mobile app must clearly communicate the outcome to the stocker so they can move on to the next item with confidence. The app displays a simple success message identifying the product and new quantity when a restock correction event succeeds, and a clear failure message when it does not, along with enough detail (e.g. invalid command, product not found, Micromart API error) for the stocker to understand what went wrong and retry if needed.

\textbf{Priority:} Essential

\textbf{Functional Requirements:}
\begin{itemize}
    \item R6.1 The system shall display a success message to the stocker, including the matched product and submitted quantity, when a restock correction event succeeds.
    \item R6.2 The system shall display a failure message to the stocker when a restock correction event fails or cannot be submitted.
    \item R6.3 The system shall distinguish between failure types shown to the stocker, including invalid command, product not found, and Micromart API error.
    \item R6.4 The system shall allow the stocker to immediately issue a new command after a result is displayed, without requiring the app to be restarted.
\end{itemize}

\section{Update Logging -- DEVIN}
For debugging and safety, the backend records every restock command it attempts to process, regardless of whether the command ultimately succeeds, fails, or is canceled by the stocker during disambiguation. Each log entry captures enough information to reconstruct what happened, including the raw input command, the parsed product name and quantity, the matched Micromart inventory identifier (if any), and the final outcome of the attempt. These logs allow the development team and the stocking company to diagnose issues after the fact without needing to reproduce them live in the field.

\textbf{Priority:} High

\textbf{Functional Requirements:}
\begin{itemize}
    \item R7.1 The system shall log every received restock command, including the raw input text and the active restocking session.
    \item R7.2 The system shall log the parsed product name and quantity extracted from each command.
    \item R7.3 The system shall log the outcome of each attempted update, including success, failure, invalid command, not found, or stocker-canceled disambiguation.
    \item R7.4 The system shall log the Micromart inventory identifier associated with a matched product when one is found.
    \item R7.5 The system shall retain logs in a form accessible to the development team for debugging purposes.
\end{itemize}

\chapter{Other Nonfunctional Requirements}

\section{Performance Requirements -- DEVIN}
Because the system is used by a stocker actively standing in front of a vending machine, responsiveness is critical to the product's value proposition of being faster than the existing Micromart website.

\begin{itemize}
    \item NFR1.1 The system shall return a result to the app within 3 seconds of receiving a parsed restock command, consistent with R2.5.
    \item NFR1.2 The system shall complete planogram fetching and index building within 10 seconds at the start of a restocking session.
    \item NFR1.3 The system shall support at least one active restocking session without degradation in response time, consistent with the single-user MVP scope.
\end{itemize}

\section{Security Requirements -- DEVIN}
Although the MVP simplifies authentication, the system must still handle credentials and inventory-affecting actions responsibly.

\begin{itemize}
    \item NFR2.1 The system shall store the manually supplied Micromart session cookie securely and shall not expose it to the mobile app beyond what is required to authenticate backend requests.
    \item NFR2.2 The system shall be designed so that the manual session cookie authentication can be replaced with a full Micromart login and two-factor authentication flow without redesigning the backend's core command-processing logic.
    \item NFR2.3 The system shall only submit restock correction events to Micromart on behalf of an authenticated, active restocking session.
    \item NFR2.4 The system shall not log the raw session cookie or other authentication credentials in the update logs described in Section 3.7.
\end{itemize}

\section{Software Quality Attributes -- DEVIN}
\begin{itemize}
    \item \textbf{Usability:} The system shall allow a stocker to issue a restock command and receive a result using only voice or short typed input, without navigating menus, so that it is faster than the existing Micromart website.
    \item \textbf{Reliability:} The system shall correctly report failures rather than silently discarding commands, and shall never submit a restock correction event for an unconfirmed or ambiguous product match.
    \item \textbf{Maintainability:} The system shall separate mobile app responsibilities (input capture, result display, disambiguation) from backend responsibilities (parsing, matching, Micromart communication) so that each can be modified independently.
    \item \textbf{Extensibility:} The system shall be structured so that support for additional vending machine platforms, such as 365, or a full login and two-factor authentication flow, can be added without a full redesign.
\end{itemize}

\section{Business Rules -- DEVIN}
\begin{itemize}
    \item BR1 The system shall support only one user and one active restocking session at a time in the initial version.
    \item BR2 The system shall not submit a restock correction event to Micromart until a command has resolved to exactly one matched product and quantity, per R4.2 and R5.1.
    \item BR3 The project team is not responsible for building or modifying the Micromart platform, the vending machine hardware, or the stocking company's existing website.
    \item BR4 The system shall focus on Micromart restocking first, with support for additional platforms, such as 365, considered out of scope for the initial version.
\end{itemize}
\end{document}
